\chapter{KẾT QUẢ THỰC NGHIỆM (ĐANG THỰC HIỆN)}
\label{ch:ket-qua}

\emph{Tính đến thời điểm báo cáo tiến độ, đồ án chưa có số liệu huấn luyện và đánh giá
chính thức. Dữ liệu đầu vào đã được chuẩn bị, quy trình xử lý đã được chạy thử qua các
bước chính, nhưng mô hình chưa được huấn luyện chính thức trên GPU. Phần này chỉ ghi
nhận trạng thái hiện tại và cố định cách báo cáo cho giai đoạn tiếp theo.}

\section{Kế hoạch thiết lập thực nghiệm}
\label{sec:kq-setup}

Thực nghiệm sẽ được chạy trong môi trường ghi nhận phiên bản hệ điều hành, Python,
PyTorch, Transformers, scikit-learn, vnstock và các thư viện dùng để tính chỉ báo kỹ
thuật. Bước tinh chỉnh PhoBERT dự kiến thực hiện trên GPU của Google Colab hoặc Kaggle
vì máy phát triển hiện không có GPU phù hợp. Tên GPU, dung lượng bộ nhớ, thời gian chạy
và checkpoint sẽ được lưu trong bản kiểm kê cấu hình.

Hạt giống ngẫu nhiên sẽ được cố định cho bước chia dữ liệu, lấy mẫu, khởi tạo mô hình và
bootstrap. Tập kiểm thử và tập giữ lại cuối cùng không được dùng để chọn siêu tham số.
Các quy tắc về mốc cắt thông tin, cách tạo nhãn và tiền xử lý chỉ học từ tập huấn luyện
đã được đặt ra trước khi chạy thực nghiệm chính.

\section{Kế hoạch đánh giá và các bảng kết quả sẽ báo cáo}
\label{sec:kq-plan}

Khi có số liệu, chương này sẽ báo cáo theo thứ tự: thống kê dữ liệu sau các cổng xử lý,
chất lượng nhãn cảm xúc và độ đồng thuận, chỉ số phân loại của PhoBERT trên tập kiểm thử,
sau đó là kết quả dự báo xu hướng theo thang mô hình.

Chỉ số chính là macro-F1, kèm balanced accuracy và precision, recall, F1 theo từng lớp
UP, FLAT và DOWN. So sánh sẽ đi từ majority, random và XGBoost chỉ giá đến LSTM chỉ
giá, LSTM hai nhánh có sentiment, rồi TFT nếu dữ liệu và thời gian cho phép. Phần đóng
góp của cảm xúc được tính bằng hiệu số giữa hai cấu hình trên cùng fold, kèm khoảng tin
cậy bootstrap khi đủ số fold.

Nguyên tắc trình bày là không điền số liệu chưa đo. Các ô còn thiếu sẽ để trạng thái
chưa điền, không dùng kết quả của công trình khác để thay thế cho dữ liệu của đề tài.
Nếu cấu hình có cảm xúc không cải thiện, kết luận phù hợp là chưa có bằng chứng về
giá trị gia tăng trong quy trình hiện tại.

\section{Dự kiến hoàn thành}
\label{sec:kq-timeline}

Bước tinh chỉnh PhoBERT và sinh xác suất cảm xúc dự kiến hoàn thành trong tuần đầu
tháng 9. Bước ghép đặc trưng, chạy baseline và LSTM hai nhánh dự kiến thực hiện trong
hai tuần tiếp theo, hướng tới hoàn thành trước hạn nộp ngày 23 tháng 9. Các mốc này là
kế hoạch, không phải kết quả đã đạt được.
